% ==============================================================================
% Chapter: Garbology — Physical Traces, Behavioral Logs, and Hermetic Scavenging
% ==============================================================================
\section{裏口の物理ストレージ：生データ監査としての考古学}

% ------------------------------------------------------------------------------
% [DRAFT CACHE / WORKING MEMORY]
% William Rathje & Cullen Murphy: "Rubbish! The Archaeology of Garbage" (1992)
% Core empirical axioms for the epistemology of physical trace analysis.
% ------------------------------------------------------------------------------
\begin{comment}
[Garbology: Core Empirical Discoveries & Zero-Trust Epistemology]

1. 「自己申告（インタビュー・メタデータ）」と「物理ログ（ゴミ箱・グラウンドトゥルース）」の完全な断絶
   - アンケートで得られる消費行動の自己申告は、回収されたゴミ袋の物的証拠と統計的に乖離する。
   - 人間は欺瞞を意図せずとも、無意識に「理想化された自己像（虚構のメタデータ）」を申告する。

2. アルコール消費量の欺瞞（過少申告のパラドックス）
   - 「ビールは週に数本」と証言した居住区のゴミ箱から、膨大な空き缶・空き瓶が発掘される。
   - 物理的痕跡（ハードウェアのログ）を突きつけられて初めて、自己の真の実行状態を認識する。

3. 「不足パニック」と食品廃棄の急増（危機時の認知エラー）
   - 1973年の牛肉不足パニック時、過剰な買い占めが発生した結果、通常時を遥かに上回る腐敗牛肉が廃棄された。
   - 欠乏への恐怖がバッファの過剰確保を生み、結果としてリソースの物理的破棄を最大化させる。

4. 「健康志向」という免罪符とジャンクフードの痕跡
   - 「野菜を多く摂取している」と申告する世帯ほど、手つかずで腐敗・廃棄された野菜の体積が大きい。
   - 実際には包装された高カロリー加工食品（スナック、簡便食）を消費し、生鮮食品の購入は心理的免罪符として機能する。

5. 埋立地の構造的実態：紙（新聞・帳票） vs プラスチックの神話崩壊
   - プラスチック汚染が象徴的に糾弾される一方、深層掘削で確認された容積の約40〜50%は紙類（新聞紙、電話帳、書類）。
   - 目立つ単一の記号をスケープゴート化し、日常を支える主幹資材の占有率を見落とす認知バイアス。

6. 「生分解」の幻想（嫌気性埋立地はタイムカプセルである）
   - 酸素と水分が遮断された深層埋立地では、有機物は容易に分解・揮発しない。
   - 数十年前の印字が鮮明に読める新聞紙、緑色を保ったレタス、皮付きのアボカドペーストが無傷で発掘される。
   - 適切な物理環境（嫌気状態）下では、生ログは半永久的に揮発せず保存される。

7. 階層ごとの廃棄プロトコルの逆転
   - 高所得者層は「環境意識」の言説を多く語るが、使い捨て包装や消費資材の物理廃棄量は最大化する。
   - 低所得者層は資材を徹底的に使い倒すため、ゴミの絶対容積および食品可食部の破棄率は極めて低い。

8. 「正面玄関（UI・言説）」をバイパスし、「裏口のゴミ箱（排出データ）」を監査せよ
   - 考古学手法を過去の遺物から「現代システムのランタイム監査」へと拡張。
   - 言説、歴史書、思想、自己申告をすべて遮断し、システムが排出した物理的ログのみから真の動作仕様を復元する。
\end{comment}

% 本文の執筆はここから継続
\noindent
人間が正面玄関で語る理念、思想、自己申告は、すべて脱臭され整えられたインターフェースに過ぎない。
本章が対象とするのは、その背後に隠された裏口のゴミ捨て場――すなわち、偽造不可能な物理ログである。

